
% !TEX root = ../main.tex
% ==============================================================================
% REFERENCIAS BIBLIOGRAFICAS Y DOCUMENTALES
% Proyecto: Campus Zero - Grupo J: SmartGrid
% Responsable: Constanza Rios
% ==============================================================================

\begin{thebibliography}{99}

% --- Antecedentes del proyecto ---
\bibitem{grupoJ2026}
Grupo J: SmartGrid,
\emph{Campus Zero: Control Automatizado de Iluminación y Climatización de Salas},
Primer Avance, asignatura Ingeniería, Informática y Sociedad,
Universidad Técnica Federico Santa María, Valparaíso, Chile,
octubre de 2026.

% --- Marco regulatorio y seguridad eléctrica ---
\bibitem{secDecreto8}
Superintendencia de Electricidad y Combustibles (SEC),
\emph{Reglamento de Seguridad de las Instalaciones de Consumo
de Energía Eléctrica, Decreto Supremo N.º 8 de 2019}.
Disponible en:
\url{https://www.sec.cl/reglamento-de-seguridad-de-las-instalaciones-de-consumo-de-energia-electrica-decreto-08/}.

\bibitem{secRIC04}
Superintendencia de Electricidad y Combustibles (SEC),
\emph{Pliego Técnico Normativo RIC N.º 04:
Conductores, Materiales y Sistemas de Canalización}.
Disponible en:
\url{https://www.sec.cl/sitio-web/wp-content/uploads/2021/01/RIC-N04-Conductores-y-Canalizaciones.pdf}.

\bibitem{secRIC05}
Superintendencia de Electricidad y Combustibles (SEC),
\emph{Pliego Técnico Normativo RIC N.º 05:
Medidas de Protección contra Tensiones Peligrosas
y Descargas Eléctricas}.
Disponible en:
\url{https://www.sec.cl/sitio-web/wp-content/uploads/2021/01/RIC-N05-Medidas-de-Proteccion-Contra-Tensiones-Peligrosas.pdf}.

\bibitem{secRIC14}
Superintendencia de Electricidad y Combustibles (SEC),
\emph{Pliego Técnico Normativo RIC N.º 14:
Exigencias de Eficiencia Energética para Edificios}.
Disponible en:
\url{https://www.sec.cl/sitio-web/wp-content/uploads/2021/01/RIC-N14-Exigencias-de-eficiencia-energetica-para-edificios.pdf}.

\bibitem{secInstaladores}
Superintendencia de Electricidad y Combustibles (SEC),
\emph{Clases de Instalador o Instaladora de Electricidad}.
Disponible en:
\url{https://www.sec.cl/area-instaladores/instaladores-electricos/}.

% --- Sensores de presencia y eficiencia energética ---
\bibitem{doeSensores}
U.S. Department of Energy,
\emph{Wireless Occupancy Sensors for Lighting Controls:
An Applications Guide for Federal Facility Managers},
2019.
Disponible en:
\url{https://www.energy.gov/sites/prod/files/2019/12/f70/wireless_occupancy_sensor_guide.pdf}.

% --- Documentación técnica de componentes ---
\bibitem{esp32Datasheet}
Espressif Systems,
\emph{ESP32 Series Datasheet},
documentación técnica oficial.
Disponible en:
\url{https://documentation.espressif.com/esp32_datasheet_en.html}.

% --- Fichas generales mencionadas en el Primer Avance ---
% Sustituir por fichas del modelo exacto cuando se confirme.

\bibitem{ledHispania}
LED Hispania Lighting,
\emph{Catálogo y fichas técnicas de paneles LED}.
Disponible en:
\url{https://www.ledhispania.es}.

\bibitem{sodecaCatalogo}
SODECA,
\emph{Catálogos técnicos de ventiladores y extractores de aire}.
Disponible en:
\url{https://www.sodeca.com}.

% --- Pendientes de la bibliografía definitiva ---
% Agregar bases oficiales del fondo institucional utilizado.
% Agregar cotizaciones específicas: producto, proveedor, fecha y URL.
% Agregar fichas de sensores PIR, relés y medidores seleccionados.
% No incluir cotizaciones que el equipo no haya consultado.

\end{thebibliography}
